\section{The initial run and aging}\label{sec:preparation}

The first run is the only part of the walk that the two fresh laws do not
describe, and Proposition~\ref{prop:first-run} says exactly how it enters. This
section turns that identity into a criterion for each zone of the range of the count.

\subsection{Statements}\label{sec:prep-statements}

The law of the initial run may depend on $N$. Throughout this section
$(\Sigma,D)=(\Sigma^{(N)},D^{(N)})$ has a law $\nu=\nu^{(N)}$ that may vary with $N$, and we
assume the tightness
\begin{equation}\label{eq:tight}
 \varepsilon_J=\sup_N\Prob\bigl(D^{(N)}>J\bigr)\to0\qquad(J\to\infty).
\end{equation}
A fixed law satisfies~\eqref{eq:tight} because $D$ is finite almost surely, and it is the case
we have in mind. The general case costs nothing in the proofs, and Section~\ref{sec:aging}
needs it. We write $\nu_\sgn$, $Q_\sgn$ and $\pi_\sgn=\Prob(\Sigma=\sgn)$ without the index $N$,
and we put
\begin{equation}\label{eq:envelope}
 \bar\nu_\sgn(n)=\sup_N\sup_{d\ge n}\nu^{(N)}_\sgn(d),\qquad\bar\nu(n)=\bar\nu_\plus(n)+\bar\nu_\minus(n),
\end{equation}
the upper envelopes of the family. For a fixed law, $\bar\nu_\plus(n)=o(T_\plus(n))$ if and only
if $\nu_\plus(n)=o(T_\plus(n))$, since $T_\plus$ is nonincreasing, and $B_n\bar\nu(n/2)\to0$ if
and only if $B_n\Prob(D=n)\to0$, since $B_n\le CB_d$ for $d\ge n/2$ by Lemma~\ref{lem:rv}(ii)
and~(v).

For a sign $\sgn$ and $b\in\Z$ define the directional defect
\begin{equation}\label{eq:defect}
 \mathcal D^\sgn_N(b)=B_N\sum_{1\le m<N/2}\nu_\sgn(N-m)\,\Pi^{\osgn}_m(b),\qquad
 \Delta_{\sgn,N}=\sup_{b\ge1}\mathcal D^\sgn_N(b),
\end{equation}
with $\Pi^{\osgn}_m(b)=0$ for $b\notin\{0,\dots,m\}$. Note that $\Pi^{\osgn}_m(0)=0$
for $m\ge1$, since a fresh walk has at least one step of its own sign. So
$\mathcal D^\sgn_N(b)=0$ for $b\le0$ and for $b\ge N/2$. Take a first run of sign
$\sgn$ that ends within the last $N/2$ steps. Times $B_N$, its contribution at the
point $k$ is $\mathcal D^\plus_N(N-k)$ for $\sgn=\plus$ and $\mathcal D^\minus_N(k)$ for
$\sgn=\minus$. So the argument $b$ is the number of steps of the opposite sign. We
say that a fixed initial run satisfies the \emph{bulk criterion} if
$B_n\Prob(D=n)\to0$.

For the big-jump zone let $L_k$ be as in~\eqref{eq:Lk} and $\widehat G_N$ as
in~\eqref{eq:Ghat}, and define for $pN<k\le N$ and $j\ge0$
\begin{align}
 \kappa_\plus(j)&=\pi_\plus u_\minus(j)\ones_{\{j\ge1\}}+\sum_{d=1}^{j}\nu_\minus(d)\,u_\minus(j-d),
 \label{eq:kappa-prep}\\
 G^\nu_N(k)&=\frac{N-L_k}{\mu_\minus}\,f_\plus(\lfloor L_k\rfloor)+T_\plus(\lfloor L_k\rfloor)\,\kappa_\plus(N-k),
 \label{eq:G-prep}\\
 \Psi_N(k)&=\sum_{L_k/2\le d<N}\nu_\plus(d)\,F^\minus_{N-d}(k-d).
 \label{eq:Psi}
\end{align}
The kernel $\kappa_\plus(j)$ is the probability that a plus run other than the
first run starts immediately after the $j$-th minus step. For a fresh start of
sign $\sgn_0$ it equals $\kappa^{(\sgn_0)}_\plus(j)$ for $j\ge1$, by~\eqref{eq:kappa}.
The kernel $G^\nu_N$ uses the limit coefficient $(N-L_k)/\mu_\minus$. The term
$\Psi_N(k)$ is the exact contribution of a first plus run that supplies at least
half of the big jump. We also use two conditions on $\nu_\plus$. The first is the
regularity condition
\begin{equation}\label{eq:Rnu}\tag{R$\nu$}
\begin{gathered}
 \sup_{d\ge n/2}\nu_\plus(d)\le C_\nu\bigl(\nu_\plus(n)+T_\plus(n)\bigr)\ \text{ for all large $n$},\\
 \lim_{\delta\downarrow0}\ \limsup_{n\to\infty}\ \sup_{|d-n|\le\delta n}
 \frac{|\nu_\plus(d)-\nu_\plus(n)|}{\nu_\plus(n)+T_\plus(n)}=0 ,
\end{gathered}
\end{equation}
where for a family both bounds are required uniformly in $N$. It says that $\nu_\plus$ is
continuous at the scale $n$ up to errors of the order of $T_\plus(n)$. Every regularly
varying $\nu_\plus$ with a negative index satisfies (R$\nu$), by Lemma~\ref{lem:rv}(i)
and~(ii), and so do $\nu_\plus=0$ and, under (LR), the fresh plus start $\nu_\plus=f_\plus$.
The second is a condition on the window averages
\[
 W_n=\frac1{B_n}\sum_{|d-n|\le B_n}\nu_\plus(d),
\]
written $W^{(N)}_n$ when the law depends on $N$.

\begin{introtheorem}[the initial run]\label{thm:preparation}
Let the law of the initial run $(\Sigma,D)$ depend on $N$ or not, subject to~\eqref{eq:tight}.
\begin{enumerate}[(a)]
\item (The exact defect.) Assume (RV) and (SPAN) for both signs. Uniformly over
$0\le k\le N$,
\begin{align*}
 B_NP_N(k)={}&h_N(x_{N,k})
 +B_NQ_\plus(N)\ones_{\{k=N\}}+B_NQ_\minus(N)\ones_{\{k=0\}}\\
 &+\mathcal D^\plus_N(N-k)+\mathcal D^\minus_N(k)+o(1).
\end{align*}
\item (The bulk criterion.) Under the hypotheses of (a),
\begin{align*}
 \sup_{1\le k<N}\bigl|B_NP_N(k)-h_N(x_{N,k})\bigr|&=\max\bigl(\Delta_{\plus,N},\Delta_{\minus,N}\bigr)+o(1),\\
 \sup_{0\le k\le N}\bigl|B_NP_N(k)-h_N(x_{N,k})\bigr|&=\max\bigl(\Delta_{\plus,N},\Delta_{\minus,N},B_NQ_\plus(N),B_NQ_\minus(N)\bigr)+o(1).
\end{align*}
The first error tends to zero if $B_N\bar\nu(N/2)\to0$ and only if $B_N\Prob(D^{(N)}=N-1)\to0$,
and the second needs in addition $B_NQ_\sgn(N)\to0$ for both signs. For a fixed law the
first error tends to zero if and only if $B_n\Prob(D=n)\to0$, and the second if and only if
$B_n\Prob(D\ge n)\to0$. For every fixed $\varepsilon>0$ the approximation holds on
$\varepsilon N\le k\le(1-\varepsilon)N$ with no condition on the initial run
beyond~\eqref{eq:tight}.
\item (The big-jump zone.) Assume (RV) and (SPAN) for both signs and (LR) for the
plus sign, and let $\lambda_N$ be as in~\eqref{eq:lambdaN}. Then
\[
 \lim_{K\to\infty}\ \limsup_{N\to\infty}\ \sup_{pN+K\lambda_NB_N\le k\le N}
 \frac{\bigl|P_N(k)-Q_\plus(N)\ones_{\{k=N\}}-G^\nu_N(k)-\Psi_N(k)\bigr|}{\widehat G_N(k)}=0 .
\]
The same holds with $\Psi_N$ removed if $\bar\nu_\plus(n)=o(T_\plus(n))$, and only if
$\nu_\plus(N-1)=o(T_\plus(N))$. For a fixed law both conditions read
$\nu_\plus(n)=o(T_\plus(n))$. If $\nu_\plus$ satisfies (R$\nu$), it holds with $\Psi_N(k)$
replaced by $\nu_\plus(\lfloor L_k\rfloor)(1+\mu_\plus u_\minus(N-k))\ones_{\{k<N\}}$ and
$\widehat G_N(k)$ replaced by $\widehat G_N(k)+\nu_\plus(\lfloor L_k\rfloor)$.
\item (The boundary layer.) Assume (RV) for the plus sign. For every fixed $j\ge0$,
\[
 P_N(N-j)=Q_\plus(N)\ones_{\{j=0\}}+T_\plus(N)\,\kappa_\plus(j)+B_N^{-1}\mathcal D^\plus_N(j)+o(T_\plus(N)),
\]
where $B_N^{-1}\mathcal D^\plus_N(j)=\sum_{1\le m<N/2}\nu_\plus(N-m)\Pi^\minus_m(j)$ does
not involve the scales. The defect term is $o(T_\plus(N))$ for every $j\ge1$ if
$\bar\nu_\plus(n)=o(T_\plus(n))$, and if it is $o(T_\plus(N))$ for one $j\ge1$, then
$\nu_\plus(N-m_0)=o(T_\plus(N))$ for some fixed $m_0=m_0(j)\ge1$. For a fixed law both
conditions read $\nu_\plus(n)=o(T_\plus(n))$, and then $P_N(N-j)/T_\plus(N)\to\kappa_\plus(j)$
for every $j\ge1$. If $\nu_\plus$ satisfies (R$\nu$), then
$B_N^{-1}\mathcal D^\plus_N(j)=\nu_\plus(N)\bigl(1+\mu_\plus u_\minus(j)\bigr)+o\bigl(\nu_\plus(N)+T_\plus(N)\bigr)$
for every $j\ge1$.
\item (Interior of the big-jump zone.) Assume the hypotheses of (c). For
$0<\varepsilon<\frac12$ let $Z^\varepsilon_N=\{k:\varepsilon N\le L_k\le(1-\varepsilon)N\}$,
and consider the property
\begin{equation}\label{eq:interior-target}
 \sup_{k\in Z^\varepsilon_N}\frac{|P_N(k)-G^\nu_N(k)|}{\widehat G_N(k)}\to0
 \quad(N\to\infty)\qquad\text{for every }\varepsilon\in(0,\tfrac12).
\end{equation}
The condition $\bar\nu_\plus(n)=o(T_\plus(n))$ implies~\eqref{eq:interior-target}.
Conversely~\eqref{eq:interior-target} implies the window condition
$W^{(2n)}_n=o(T_\plus(n))$. If (LR) holds for the minus sign as well, then
$\sup_{\varepsilon N/4\le n\le N}W^{(N)}_n=o(T_\plus(N))$ for every $\varepsilon\in(0,\frac12)$
implies~\eqref{eq:interior-target}. For a fixed law the last two conditions read
$W_n=o(T_\plus(n))$.
\end{enumerate}
At the plus end the criteria are nested. Since $\nu_\plus\le Q_\plus$, the atom
condition $\sup_NQ_\plus(n)=o(T_\plus(n))$ implies the big-jump condition
$\bar\nu_\plus(n)=o(T_\plus(n))$. Since $B_nT_\plus(n)\to0$, the big-jump condition
implies $B_n\bar\nu_\plus(n/2)\to0$, which is the plus half of the bulk criterion.
\end{introtheorem}

The zone in (c) starts where the zone of Theorem~\ref{thm:bigjump}(b) starts, and by
Remark~\ref{rem:theta} its factor $\lambda_N$ cannot be replaced by a constant, already
for a fresh start. For fresh starts Theorem~\ref{thm:window} gives the sharp factor. At $k=N$ we have $\Psi_N(N)=0$, $G^\nu_N(N)=0$ and
$\mathcal D^\plus_N(0)=0$, so (c) and (d) both reduce to the identity
$P_N(N)=Q_\plus(N)$ of Proposition~\ref{prop:first-run}. Part (e) refines (c) on
the interior of the zone, where~\eqref{eq:interior-target} is the statement of (c)
without $\Psi_N$. The regular case of (c) gives one formula on the whole zone up to the
end, with the term $\nu_\plus(\lfloor L_k\rfloor)(1+\mu_\plus u_\minus(N-k))$ for the first
run as the giant run, and the regular case of (d) is its value at $j=N-k$ fixed, since
$\nu_\plus(\lfloor L_k\rfloor)=\nu_\plus(N)(1+o(1))+o(T_\plus(N))$ there by (R$\nu$).

\begin{remark}[the prepared boundary kernel]\label{rem:kappa-prep}
By~\eqref{eq:renewal-thm} and dominated convergence, $\kappa_\plus(j)\to1/\mu_\minus$,
the fresh value. With generating functions $F$, $U=1/(1-F)$ and $\widehat\nu_\minus$ of
$f_\minus$, $u_\minus$ and $\nu_\minus$, the identity $\kappa_\plus=u_\minus$ on $j\ge1$
reads $\widehat\nu_\minus U=\pi_\minus(U-1)$, that is, $\nu_\minus=\pi_\minus f_\minus$. So
the prepared kernel is the fresh one exactly when the first run, given that it is a
minus run, has the law of $X_\minus$.
\end{remark}

\subsection{Tools}\label{sec:prep-tools}

For a horizon $n$ and a point $k>pn$ we write $L^{(n)}_k=(k-pn)/(1-p)$ for the
quantity~\eqref{eq:Lk} at the horizon $n$, and $\widehat G_n(k)$, $G^{(\sgn_0)}_n(k)$
and $\lambda_n$ for~\eqref{eq:Ghat}, \eqref{eq:G} and~\eqref{eq:lambdaN} at the
horizon $n$. Recall from Lemma~\ref{lem:rv}(v) that $B_n$ is nondecreasing and
regularly varying with index $1/\alpha_*$.

Fix $c\in(0,1)$. We have $\lambda_n=1+\log\max(1,r(n))$ with
$r(n)=b_{\minus,n}/b_{\plus,n}$, which is regularly varying with index
$\varrho=1/\alpha_\minus-1/\alpha_\plus$. By the uniform convergence
theorem~\cite[Theorem~1.5.2]{bgt1989}, $r(n)/r(N)=(n/N)^{\varrho}(1+o(1))$ uniformly in
$n\in[cN,N]$. If $\varrho\ge0$, then $(n/N)^\varrho\le1$, so $r(n)\le r(N)(1+o(1))$.
If $\varrho<0$, then $r(N)\to0$, so $r(n)\le c^{\varrho}r(N)(1+o(1))<1$ for large $N$ and
$\lambda_n=\lambda_N=1$. Since $x\mapsto\log\max(1,x)$ is nondecreasing and changes by
at most $\log(1+\delta)$ when $x$ is multiplied by $1+\delta$, in both cases
$\lambda_n\le\lambda_N+o(1)$ uniformly in $n\in[cN,N]$. Since $\lambda_N\ge1$,
\begin{equation}\label{eq:lambda-unif}
 \lambda_n\le2\lambda_N\qquad(cN\le n\le N)
\end{equation}
once $N$ is large. By Lemma~\ref{lem:rv}(v), $\lambda_NB_N=o(N)$. So for each $c>0$
the points with $k-pN\ge cN$ lie in the zone of Theorem~\ref{thm:preparation}(c) for
every $K$ once $N$ is large.

\begin{lemma}[tightness of the count]\label{lem:prep-tight}
Assume (RV) and (SPAN). For a fresh start of either sign,
$\sup_{n\ge1}\Prob(|A_n-pn|>RB_n)\to0$ as $R\to\infty$.
\end{lemma}

\begin{proof}
Fix the starting sign and put $\varepsilon_n=\sup_k|B_nF_n(k)-h_n(x_{n,k})|$, which
tends to zero by Theorem~\ref{thm:bulk}(a). Let $R\ge1$ and let $k$ run over the
integers with $|k-pn|\le RB_n$. The points $x_{n,k}$ are spaced $1/B_n$ apart, and the
intervals $[x_{n,k},x_{n,k}+1/B_n)$ cover $[-R+1/B_n,R]$. On each of them $h_n$
exceeds $h_n(x_{n,k})$ by at most $C_2/B_n$, by Lemma~\ref{lem:profile}(a). There are
at most $2RB_n+1$ such $k$. Hence, for $B_n\ge1$,
\begin{align*}
 \Prob(|A_n-pn|\le RB_n)&\ge\sum_k\frac{h_n(x_{n,k})-\varepsilon_n}{B_n}
 \ge\int_{-R+1/B_n}^{R}h_n-3R\Bigl(\varepsilon_n+\frac{C_2}{B_n}\Bigr)\\
 &\ge1-\varepsilon'_h(R-1)-3R\Bigl(\varepsilon_n+\frac{C_2}{B_n}\Bigr).
\end{align*}
So $\limsup_n\Prob(|A_n-pn|>RB_n)\le\varepsilon'_h(R-1)$, which tends to zero by
Lemma~\ref{lem:profile}(b). For each of the finitely many remaining $n$ the
probability is zero once $RB_n>n$.
\end{proof}

The measure $m\mapsto\Pi^\minus_m(j)$ of the Green identity sits within $O(B_N)$ of
$j/(1-p)$, the time at which a typical walk collects $j$ minus steps.

\begin{lemma}[the Green measure is concentrated]\label{lem:green-conc}
Assume (RV) and (SPAN). For every $\varepsilon>0$ there are $R$ and $N_0$ such that
for all $N\ge N_0$ and all $1\le j\le N$,
\[
 \sum_{m\ge1:\ |m-j/(1-p)|>RB_N}\Pi^\minus_m(j)\le\varepsilon .
\]
\end{lemma}

\begin{proof}
Take a fresh minus walk continued forever, as in the proof of
Lemma~\ref{lem:green}. Its minus runs $X^\minus_1,X^\minus_2,\dots$ and its plus runs
$X^\plus_1,X^\plus_2,\dots$ are independent, and the $i$-th plus run comes right after
the $i$-th minus run. Let $C_m$ be the number of minus steps among the first $m$, let
$t_j$ be the step of the $j$-th minus step, and let $r$ be the index of the minus run
that contains it. Put $\xi_j=X^\plus_r$ if $S^\minus_r=j$ and $\xi_j=0$ otherwise. As
in the proof of Lemma~\ref{lem:green}, $C_m=j$ exactly for $t_j\le m\le t_j+\xi_j$.
The index $r$, the event $\{S^\minus_r=j\}$ and $t_j=j+X^\plus_1+\dots+X^\plus_{r-1}$
are functions of the minus runs and of $X^\plus_1,\dots,X^\plus_{r-1}$. So $X^\plus_r$
is independent of them and has the law of $X_\plus$.

Put $c=j/(1-p)$, $\mathcal S=\{m:|m-c|>RB_N\}$ and $\mathcal V=\{m:|m-c|\le RB_N/2\}$.
If $t_j\in\mathcal V$ and $\xi_j<RB_N/2$, then the steps $t_j,\dots,t_j+\xi_j$ all lie
outside $\mathcal S$. Hence, by Tonelli's theorem,
\[
 \sum_{m\in\mathcal S}\Pi^\minus_m(j)=\E\,\#\{m\in\mathcal S:C_m=j\}
 \le\E\bigl[(1+X^\plus_r)\ones_{\{t_j\notin\mathcal V\}}\bigr]+\E\bigl[(1+X^\plus_r)\ones_{\{X^\plus_r\ge RB_N/2\}}\bigr].
\]
By the independence just noted, the first term is
$(1+\mu_\plus)\Prob(t_j\notin\mathcal V)$ and the second is
$\E[(1+X_\plus)\ones_{\{X_\plus\ge RB_N/2\}}]$, which tends to zero as $N\to\infty$
because $\mu_\plus<\infty$.

It remains to bound $\Prob(t_j\notin\mathcal V)$. Let $m_1=\lfloor c+RB_N/2\rfloor$
and $m_2=\lceil c-RB_N/2\rceil$. If $t_j\notin\mathcal V$, then $t_j>m_1$ or
$t_j\le m_2$. Now $\{t_j>m_1\}=\{C_{m_1}<j\}=\{A_{m_1}>m_1-j\}$, and
$m_1-j-pm_1=(1-p)m_1-j\ge(1-p)(RB_N/2-1)\ge(1-p)RB_N/4$ for large $N$. In the same
way $\{t_j\le m_2\}=\{A_{m_2}\le m_2-j\}$ and $(1-p)m_2-j\le-(1-p)RB_N/4$, and this
event is empty if $m_2<1$. For $N\ge N_0$ we have $RB_N\le N$, so $m_1,m_2\le c_9N$
with $c_9=2+1/(1-p)$, and $B_{m_i}\le B_{c_9N}\le c_{10}B_N$ for some $c_{10}$ by
Lemma~\ref{lem:rv}(ii). Hence
\[
 \Prob(t_j\notin\mathcal V)\le2\sup_{n\ge1}\Prob\Bigl(|A_n-pn|>\frac{(1-p)R}{4c_{10}}\,B_n\Bigr),
\]
which is small for large $R$ by Lemma~\ref{lem:prep-tight}.
\end{proof}

Finally we record the fresh big-jump asymptotics at a general horizon. For a
horizon $n$ and a point $k\le n$ with $L=L^{(n)}_k\ge1$ let
$\overline G^{(\sgn_0)}_n(k)=(n-L)f_\plus(\lfloor L\rfloor)/\mu_\minus
+T_\plus(\lfloor L\rfloor)\kappa^{(\sgn_0)}_\plus(n-k)$ be the kernel~\eqref{eq:G} at the
horizon $n$ with the limit coefficient. Assume (RV) and (SPAN) for both signs and
(LR) for the plus sign. By the last clause of Theorem~\ref{thm:bigjump}(b), applied at
the horizon $n$, for every $\varepsilon>0$ there are $K_0$ and $n_0$ such that
\begin{equation}\label{eq:unif-bigjump}
 \bigl|F^{\sgn_0}_n(k)-\overline G^{(\sgn_0)}_n(k)\bigr|\le\varepsilon\,\widehat G_n(k)
 \qquad(n\ge n_0,\ pn+K_0\lambda_nB_n\le k\le n)
\end{equation}
for both starting signs, and by~\eqref{eq:G-size} both
$\overline G^{(\sgn_0)}_n(k)$ and $G^{(\sgn_0)}_n(k)$ are at most $C_{\mathrm G}\widehat G_n(k)$
whenever $L^{(n)}_k\ge1$.

\subsection{The bulk, parts (a) and (b)}\label{sec:prep-bulk}

\begin{proof}[Proof of Theorem~\ref{thm:preparation}(a)]
By the definition of $\Pi^\sgn_m$ we have $F^\minus_{N-d}(k-d)=\Pi^\minus_{N-d}(N-k)$ and
$F^\plus_{N-d}(k)=\Pi^\plus_{N-d}(k)$. Fix $J\ge1$. By
Proposition~\ref{prop:first-run}, for $N>2J$,
\[
 B_NP_N(k)=B_NQ_\plus(N)\ones_{\{k=N\}}+B_NQ_\minus(N)\ones_{\{k=0\}}+E_1(k)+E_2(k)+E_3(k),
\]
where $E_1$, $E_2$ and $E_3$ collect the first-run lengths $d\le J$, $J<d\le N/2$ and
$N/2<d<N$, each multiplied by $B_N$. In $E_3$ put $m=N-d$, so that $1\le m<N/2$.
By~\eqref{eq:defect} this gives exactly $E_3(k)=\mathcal D^\plus_N(N-k)+\mathcal D^\minus_N(k)$.

For $E_2$ the horizon is $N-d\ge N/2$. By Theorem~\ref{thm:bulk}(b) there is $C_{10}$
with $B_nF^{\sgn_0}_n(k)\le C_{10}$ for all $n$, $k$ and $\sgn_0$. By
Lemma~\ref{lem:rv}(v), $B_N/B_{N-d}\le B_N/B_{\lceil N/2\rceil}\le C$. Hence
$0\le E_2(k)\le CC_{10}\varepsilon_J$ for all $k$, by~\eqref{eq:tight}.

For $E_1$ we claim that for each fixed $d\ge1$
\begin{equation}\label{eq:prep-shift}
 \sup_{k\in\Z}\bigl|B_NF^\minus_{N-d}(k-d)-h_N(x_{N,k})\bigr|\to0,\qquad
 \sup_{k\in\Z}\bigl|B_NF^\plus_{N-d}(k)-h_N(x_{N,k})\bigr|\to0 .
\end{equation}
We prove the first, and the second is the same with $pd$ in place of $-(1-p)d$
below. Put $n=N-d$. By Theorem~\ref{thm:bulk}(a),
$B_nF^\minus_n(k-d)=h_n(y_k)+o(1)$ uniformly in $k$, with
$y_k=(k-d-pn)/B_n=(k-pN-(1-p)d)/B_n$. Also $B_N/B_n\to1$ and $h_n$ is bounded. By
Lemma~\ref{lem:profile}(c), $\lVert h_n-h_N\rVert_\infty\to0$. Fix $R$. If
$|x_{N,k}|\le R$, then $|y_k-x_{N,k}|\le R|B_N/B_n-1|+(1-p)d/B_n\to0$, and the
Lipschitz bound of Lemma~\ref{lem:profile}(a) gives $|h_N(y_k)-h_N(x_{N,k})|\to0$. If
$|x_{N,k}|>R$, then $|y_k|\ge R/2$ for large $N$, and both values are at most
$\varepsilon_h(R/2)$. Let $N\to\infty$ and then $R\to\infty$. This
proves~\eqref{eq:prep-shift}. Summing it over the finitely many $d\le J$ gives
$\sup_k|E_1(k)-\Prob(D\le J)h_N(x_{N,k})|\to0$.

Let $e_N(k)$ be the difference between $B_NP_N(k)$ and the right side of (a)
without the term $o(1)$. Collecting the three parts,
$\limsup_{N\to\infty}\sup_{0\le k\le N}|e_N(k)|\le(C_1+CC_{10})\varepsilon_J$, with $C_1$
the bound of Lemma~\ref{lem:profile}(a). This tends to zero as $J\to\infty$,
by~\eqref{eq:tight}.
\end{proof}

\begin{proof}[Proof of Theorem~\ref{thm:preparation}(b)]
The two defects are nonnegative. If $\mathcal D^\plus_N(N-k)\ne0$, then some $m<N/2$ has
$\Pi^\minus_m(N-k)>0$, so $N-k\le m<N/2$ and $k>N/2$. In the same way
$\mathcal D^\minus_N(k)\ne0$ forces $k<N/2$. So the two defects have disjoint supports
in $1\le k<N$, and as $k$ runs over $1\le k<N$ the arguments $N-k$ and $k$ run over
all $b$ at which $\mathcal D^\sgn_N(b)$ can be positive. Hence
$\sup_{1\le k<N}\bigl(\mathcal D^\plus_N(N-k)+\mathcal D^\minus_N(k)\bigr)=\max(\Delta_{\plus,N},\Delta_{\minus,N})$.
Part (a) and the inequality $|\sup|f+g|-\sup|f||\le\sup|g|$ then give
\[
 \sup_{1\le k<N}\bigl|B_NP_N(k)-h_N(x_{N,k})\bigr|=\max(\Delta_{\plus,N},\Delta_{\minus,N})+o(1).
\]
At $k=N$ part (a) gives $B_NP_N(N)=B_NQ_\plus(N)+h_N(x_{N,N})+o(1)$, since
$\mathcal D^\plus_N(0)=0$ and $\mathcal D^\minus_N(N)=0$. Here $x_{N,N}=(1-p)N/B_N\to\infty$,
so $h_N(x_{N,N})\to0$ by Lemma~\ref{lem:profile}(b). The same holds at $k=0$ with
$Q_\minus$. This proves the formula for the error over the whole range.

Now we compare the defects with the atoms. The term $m=1$ of~\eqref{eq:defect} at
$b=1$ is $B_N\nu_\plus(N-1)\Pi^\minus_1(1)=B_N\nu_\plus(N-1)$, since a fresh minus walk
of one step has one minus step. So $\Delta_{\plus,N}\ge B_N\nu_\plus(N-1)$. In the
other direction, $\nu_\plus(N-m)\le\bar\nu_\plus(N/2)$ for $m<N/2$, and the Green
identity, Lemma~\ref{lem:green}, together with $u_\minus\le1$ gives
\begin{equation}\label{eq:defect-upper}
 \Delta_{\plus,N}\le(1+\mu_\plus)\,B_N\sup_{d>N/2}\nu_\plus(d)\le(1+\mu_\plus)\,B_N\,\bar\nu_\plus(N/2).
\end{equation}
The same bounds hold for the minus sign. So $\max(\Delta_{\plus,N},\Delta_{\minus,N})\to0$ if
$B_N\bar\nu(N/2)\to0$ and only if $B_N\Prob(D^{(N)}=N-1)\to0$, and for a fixed law both
conditions read $B_n\Prob(D=n)\to0$, by the remark after~\eqref{eq:envelope} and
$B_N/B_{N-1}\to1$. This is the interior criterion. For the whole range we need in
addition $B_NQ_\sgn(N)\to0$ for both signs, and for a fixed law, since $\nu_\sgn\le Q_\sgn$,
the two conditions together read $B_n\Prob(D\ge n)\to0$.

Finally let $\varepsilon N\le k\le(1-\varepsilon)N$. A nonzero term of
$\mathcal D^\plus_N(N-k)$ has $m\ge N-k\ge\varepsilon N$, and Theorem~\ref{thm:bulk}(b)
gives $\Pi^\minus_m(N-k)\le C_{10}/B_m\le C_{10}/B_{\lceil\varepsilon N\rceil}$. Hence
$\mathcal D^\plus_N(N-k)\le C_{10}(B_N/B_{\lceil\varepsilon N\rceil})\,Q_\plus(\lceil N/2\rceil)\to0$,
since $Q_\plus(\lceil N/2\rceil)\le\varepsilon_J$ once $N>2J$. The minus defect is the same, and part (a)
gives the approximation on this range.
\end{proof}

\subsection{The boundary layer, part (d)}\label{sec:prep-boundary}

\begin{proof}[Proof of Theorem~\ref{thm:preparation}(d)]
The proof uses only (RV) for the plus sign, through Theorem~\ref{thm:bigjump}(a) and
the bound $T_\plus(\lceil N/2\rceil)\le CT_\plus(N)$ of Lemma~\ref{lem:rv}(ii). Fix
$j\ge0$ and let $N>2j$. Proposition~\ref{prop:first-run} at $k=N-j$ gives
\[
 P_N(N-j)=Q_\plus(N)\ones_{\{j=0\}}+\sum_{d<N}\nu_\minus(d)\,F^\plus_{N-d}(N-j)
 +\sum_{d<N}\nu_\plus(d)\,F^\minus_{N-d}(N-d-j).
\]
In the first sum the fresh plus walk on $N-d$ steps has $N-j$ plus steps, so
$d\le j$, and it then has $j-d$ minus steps. For each such $d$,
Theorem~\ref{thm:bigjump}(a) at the horizon $N-d$ and~\eqref{eq:kappa} give
$F^\plus_{N-d}(N-j)/T_\plus(N-d)\to u_\minus(j-d)$, and $T_\plus(N-d)/T_\plus(N)\to1$ by
(RV). The sum has at most $j$ terms. So the first sum equals
$T_\plus(N)\bigl(\sum_{d=1}^j\nu_\minus(d)u_\minus(j-d)+o(1)\bigr)$.

Split the second sum at $J$ and at $N/2$, with $J$ fixed. For $d\le J$,
Theorem~\ref{thm:bigjump}(a) for a fresh minus start and~\eqref{eq:kappa} give
$F^\minus_{N-d}(N-d-j)/T_\plus(N-d)\to u_\minus(j)\ones_{\{j\ge1\}}$, so this part equals
$T_\plus(N)\bigl(\Prob(\Sigma=\plus,D\le J)u_\minus(j)\ones_{\{j\ge1\}}+o(1)\bigr)$. For
$J<d\le N/2$ the horizon is $n=N-d\ge N/2$. The convergent sequence
$F^\minus_n(n-j)/T_\plus(n)$ is bounded by some $C_j$ for $n\ge N/2$ and large $N$,
and $T_\plus(n)\le T_\plus(\lceil N/2\rceil)\le CT_\plus(N)$. So this part is at
most $CC_j\varepsilon_JT_\plus(N)$, by~\eqref{eq:tight}. For $N/2<d<N$ put $m=N-d$. Then
$F^\minus_{N-d}(N-d-j)=\Pi^\minus_m(j)$, and by~\eqref{eq:defect} this part equals
$B_N^{-1}\mathcal D^\plus_N(j)$ exactly. Collecting terms and using~\eqref{eq:kappa-prep},
\[
 \limsup_{N\to\infty}\frac{\bigl|P_N(N-j)-Q_\plus(N)\ones_{\{j=0\}}-T_\plus(N)\kappa_\plus(j)
 -B_N^{-1}\mathcal D^\plus_N(j)\bigr|}{T_\plus(N)}\le(1+CC_j)\,\varepsilon_J,
\]
and the right side tends to zero as $J\to\infty$. This is the expansion.

Let $j\ge1$. If $\bar\nu_\plus(n)=o(T_\plus(n))$, then the Green identity,
Lemma~\ref{lem:green}, and $u_\minus\le1$ give
$B_N^{-1}\mathcal D^\plus_N(j)\le(1+\mu_\plus)\bar\nu_\plus(N/2)
=o(T_\plus(\lceil N/2\rceil))=o(T_\plus(N))$. Conversely, suppose that the defect term is
$o(T_\plus(N))$ for one $j\ge1$. By Lemma~\ref{lem:green}
the weights $\Pi^\minus_m(j)$, $m\ge1$, add up to $1+\mu_\plus u_\minus(j)\ge1$, so
there is an $m_0=m_0(j)\ge1$ with $\Pi^\minus_{m_0}(j)>0$, and it does not depend on
$N$. (If $T_\minus(j)>0$ one may take $m_0=j$, with $\Pi^\minus_j(j)=T_\minus(j)$.)
For $N>2m_0$ the term $m=m_0$ of~\eqref{eq:defect} gives
$B_N^{-1}\mathcal D^\plus_N(j)\ge\nu_\plus(N-m_0)\Pi^\minus_{m_0}(j)$, so
$\nu_\plus(N-m_0)=o(T_\plus(N))$. For a fixed law, $T_\plus(N)/T_\plus(N-m_0)\to1$ gives
$\nu_\plus(n)=o(T_\plus(n))$, and then the expansion gives $P_N(N-j)/T_\plus(N)\to\kappa_\plus(j)$.

Finally let $\nu_\plus$ satisfy (R$\nu$) and $j\ge1$, and put $r_N=\nu_\plus(N)+T_\plus(N)$.
Fix $M$. For $m\le M$ the second part of (R$\nu$) gives
$|\nu_\plus(N-m)-\nu_\plus(N)|\le\epsilon_Nr_N$ with $\epsilon_N\to0$, and for $M<m<N/2$ the first
part gives $\nu_\plus(N-m)+\nu_\plus(N)\le(C_\nu+1)r_N$. By Lemma~\ref{lem:green} the weights
$\Pi^\minus_m(j)$ have the finite sum $1+\mu_\plus u_\minus(j)$. Hence
\begin{multline*}
 \Bigl|B_N^{-1}\mathcal D^\plus_N(j)-\nu_\plus(N)\bigl(1+\mu_\plus u_\minus(j)\bigr)\Bigr|\\
 \le\epsilon_N(1+\mu_\plus)r_N+(C_\nu+1)r_N\sum_{m>M}\Pi^\minus_m(j)+\nu_\plus(N)\sum_{m\ge N/2}\Pi^\minus_m(j),
\end{multline*}
and letting $N\to\infty$, then $M\to\infty$, the right side is $o(r_N)$.
\end{proof}

\subsection{The big-jump zone, parts (c) and (e)}\label{sec:prep-bigjump}

\begin{proof}[Proof of the decomposition in Theorem~\ref{thm:preparation}(c)]
Fix $\varepsilon\in(0,1]$ and take $K_0$ and $n_0$ from~\eqref{eq:unif-bigjump}. Let
$K\ge4K_0$ and $J\ge1$. Let $pN+K\lambda_NB_N\le k\le N$, and write $L=L_k$ and
$j=N-k$. Then $L\ge K\lambda_NB_N/(1-p)\ge KB_N/(1-p)$, since $\lambda_N\ge1$, and
$N-L=j/(1-p)$. Take $N$ so large that $\min(p,\frac12)N\ge n_0$, that $J<L/2$, and
that~\eqref{eq:lambda-unif} holds with $c=\min(p,\frac12)$.
By Proposition~\ref{prop:first-run},
\[
 P_N(k)=Q_\plus(N)\ones_{\{k=N\}}+E_1+E_2+E_3+E_4+\Psi_N(k),
\]
where $E_1$ and $E_2$ collect $\nu_\minus(d)F^\plus_{N-d}(k)$ over $d\le\min(J,j)$ and
over $J<d\le j$, and $E_3$ and $E_4$ collect $\nu_\plus(d)F^\minus_{N-d}(k-d)$ over
$d\le J$ and over $J<d<L/2$. The terms $\nu_\minus(d)F^\plus_{N-d}(k)$ with $d>j$
vanish, since then $k>N-d$.

\emph{Step 1 (geometry).} Let $d\le j$ and $n=N-d$. Then $n\ge k>pN$, so
$n\ge n_0$, and $L^{(n)}_k=L+pd/(1-p)\ge L$. Also
$k-pn\ge k-pN\ge K\lambda_NB_N\ge(K/2)\lambda_nB_n$, by~\eqref{eq:lambda-unif} and
since $B$ is nondecreasing. So $k$ lies in the zone of~\eqref{eq:unif-bigjump} at the
horizon $n$, and $\widehat G_n(k)\le(N/L)T_\plus(\lfloor L\rfloor)=\widehat G_N(k)$
because $T_\plus$ is nonincreasing. Let $d<L/2$ and $n=N-d\ge N/2$. Then
$L^{(n)}_{k-d}=L-d\ge L/2$, and
$(k-d)-pn=(1-p)(L-d)\ge K\lambda_NB_N/2\ge(K/4)\lambda_nB_n$, again
by~\eqref{eq:lambda-unif}, so $k-d$ lies in that zone at the horizon $n$. By
Lemma~\ref{lem:rv}(ii), $\widehat G_n(k-d)\le(2N/L)T_\plus(\lfloor L/2\rfloor)\le
C_*\widehat G_N(k)$ for large $N$.

\emph{Step 2 (intermediate first-run lengths).} By Step 1
and~\eqref{eq:unif-bigjump}, every fresh probability in $E_2$ and $E_4$ is at most
$(C_{\mathrm G}+1)$ times the corresponding $\widehat G_n$, since $\varepsilon\le1$.
Hence $E_2+E_4\le(C_{\mathrm G}+1)(1+C_*)\,\varepsilon_J\,\widehat G_N(k)$, by~\eqref{eq:tight}.

\emph{Step 3 (short first-run lengths).} Let $d\le J$. By Step 1
and~\eqref{eq:unif-bigjump},
\[
 |F^\plus_{N-d}(k)-\overline G^{(\plus)}_{N-d}(k)|\le\varepsilon\widehat G_N(k),\qquad
 |F^\minus_{N-d}(k-d)-\overline G^{(\minus)}_{N-d}(k-d)|\le\varepsilon C_*\widehat G_N(k).
\]
With $L'=L+pd/(1-p)$ and $(N-d)-L'=(j-d)/(1-p)$, and with~\eqref{eq:kappa}, the
kernels are
\begin{align*}
 \overline G^{(\plus)}_{N-d}(k)&=\frac{j-d}{(1-p)\mu_\minus}f_\plus(\lfloor L'\rfloor)+T_\plus(\lfloor L'\rfloor)\,u_\minus(j-d),\\
 \overline G^{(\minus)}_{N-d}(k-d)&=\frac{N-L}{\mu_\minus}f_\plus(\lfloor L-d\rfloor)+T_\plus(\lfloor L-d\rfloor)\,u_\minus(j)\ones_{\{j\ge1\}} .
\end{align*}
By (LR) and (RV), $f_\plus(m+i)/f_\plus(m)\to1$ and $T_\plus(m+i)/T_\plus(m)\to1$ as
$m\to\infty$, uniformly in $|i|\le2J/(1-p)$. Let $r_N$ be the largest deviation
from one of these ratios over $m\ge KB_N/(2(1-p))$. Then $r_N\to0$. Note that
$(j-d)/(1-p)=(N-L)-d/(1-p)$ and that $f_\plus(\lfloor L'\rfloor)\le2CT_\plus(\lfloor L\rfloor)/L$
for large $N$, by (LR). Hence
\[
 \Bigl|\overline G^{(\plus)}_{N-d}(k)-\frac{N-L}{\mu_\minus}f_\plus(\lfloor L\rfloor)-T_\plus(\lfloor L\rfloor)u_\minus(j-d)\Bigr|
 \le C\Bigl(r_N+\frac JL\Bigr)\widehat G_N(k),
\]
and $\overline G^{(\minus)}_{N-d}(k-d)$ differs from
$(N-L)f_\plus(\lfloor L\rfloor)/\mu_\minus+T_\plus(\lfloor L\rfloor)u_\minus(j)\ones_{\{j\ge1\}}$ by at
most $Cr_N\widehat G_N(k)$. Here we used $(N-L)f_\plus(\lfloor L\rfloor)\le C\widehat G_N(k)$
and $T_\plus(\lfloor L\rfloor)\le\widehat G_N(k)$.

Now we sum over $d$. In $E_3$ the weights sum to $\Prob(\Sigma=\plus,D\le J)$, which
differs from $\pi_\plus$ by at most $\varepsilon_J$. In $E_1$ the weights of the interior
term sum to $\Prob(\Sigma=\minus,D\le\min(J,j))$. If $j\ge J$ this differs from
$\pi_\minus$ by at most $\varepsilon_J$. If $j<J$, then $N-L<J/(1-p)$ and the whole
interior term is at most $CJ\,T_\plus(\lfloor L\rfloor)/L\le C(J/L)\widehat G_N(k)$. The
boundary terms in $E_1$ sum to $T_\plus(\lfloor L\rfloor)\sum_{d\le\min(J,j)}\nu_\minus(d)u_\minus(j-d)$,
which differs from $T_\plus(\lfloor L\rfloor)\sum_{d\le j}\nu_\minus(d)u_\minus(j-d)$ by at
most $\varepsilon_J\widehat G_N(k)$, since $u_\minus\le1$. By~\eqref{eq:kappa-prep}
and~\eqref{eq:G-prep},
\[
 |E_1+E_3-G^\nu_N(k)|\le C\Bigl(\varepsilon+r_N+\frac JL+\varepsilon_J\Bigr)\widehat G_N(k).
\]

\emph{Step 4 (conclusion).} Steps 2 and 3 and $J/L\le(1-p)J/(KB_N)$ give, for
$K\ge4K_0$ and large $N$,
\begin{multline*}
 \sup_{pN+K\lambda_NB_N\le k\le N}\frac{|P_N(k)-Q_\plus(N)\ones_{\{k=N\}}-G^\nu_N(k)-\Psi_N(k)|}{\widehat G_N(k)}\\
 \le C\Bigl(\varepsilon+r_N+\frac{(1-p)J}{KB_N}+\varepsilon_J\Bigr).
\end{multline*}
Let $N\to\infty$, then $J\to\infty$. The bound becomes $C\varepsilon$ for every
$K\ge4K_0(\varepsilon)$, and $\varepsilon$ is arbitrary.
\end{proof}

\begin{proof}[Proof of the criteria in Theorem~\ref{thm:preparation}(c)]
\emph{The whole zone.} By the decomposition, the statement without $\Psi_N$ holds if
and only if $\lim_K\limsup_N\sup_k\Psi_N(k)/\widehat G_N(k)=0$, the supremum being over
the zone. Suppose $\bar\nu_\plus(n)=o(T_\plus(n))$. Every $d$ in~\eqref{eq:Psi} has
$d\ge L_k/2$, and Lemma~\ref{lem:green} with $u_\minus\le1$ gives, with $j=N-k$,
\begin{equation}\label{eq:Psi-upper}
 \Psi_N(k)\le\sup_{d\ge L_k/2}\nu_\plus(d)\sum_{m\ge1}\Pi^\minus_m(j)
 \le(1+\mu_\plus)\,\bar\nu_\plus(L_k/2),
\end{equation}
because $F^\minus_{N-d}(k-d)=\Pi^\minus_{N-d}(j)$ and distinct $d$ give distinct
$m=N-d$. The right side is $o(T_\plus(\lceil L_k/2\rceil))=o(T_\plus(\lfloor L_k\rfloor))$
by Lemma~\ref{lem:rv}(ii), and $T_\plus(\lfloor L_k\rfloor)\le\widehat G_N(k)$.
Conversely, for large $N$ the point $k=N-1$ lies in the zone for every $K$, because
$\lambda_NB_N=o(N)$, and it has $L=N-1/(1-p)\ge N/2$. The term $d=N-1$
of~\eqref{eq:Psi} gives $\Psi_N(N-1)\ge\nu_\plus(N-1)F^\minus_1(0)=\nu_\plus(N-1)$, and
$\widehat G_N(N-1)\le2T_\plus(\lfloor N/2\rfloor)\le CT_\plus(N)$. So the criterion
forces $\nu_\plus(N-1)=o(T_\plus(N))$, which for a fixed law is $\nu_\plus(n)=o(T_\plus(n))$.

\emph{The regular case.} Let $\nu_\plus$ satisfy (R$\nu$) and fix $\varepsilon>0$. Take $R$
and $N_0$ from Lemma~\ref{lem:green-conc}, and $\delta\in(0,\frac12]$ and $n_\delta$ such that the first
part of (R$\nu$) holds for $n\ge n_\delta$ and
$|\nu_\plus(d)-\nu_\plus(n)|\le\varepsilon(\nu_\plus(n)+T_\plus(n))$ for $|d-n|\le\delta n$ and
$n\ge n_\delta$. Let $K\ge R(1-p)/\delta+2$ and let $k<N$ lie in the zone, so $j=N-k\ge1$ and
$L_k\ge KB_N/(1-p)$, and put $n=\lfloor L_k\rfloor$ and $r=\nu_\plus(n)+T_\plus(n)$. For large $N$
we have $n\ge n_\delta$. Let $\mathcal V=\{d\le N-1:|d-L_k|\le RB_N\}$. Since $m=N-d$ satisfies
$|m-j/(1-p)|=|d-L_k|$, Lemma~\ref{lem:green-conc} gives
$\sum_{d\le N-1,\,d\notin\mathcal V}\Pi^\minus_{N-d}(j)\le\varepsilon$, and Lemma~\ref{lem:green}
then gives $\sum_{d\in\mathcal V}\Pi^\minus_{N-d}(j)\in[1+\mu_\plus u_\minus(j)-\varepsilon,\,1+\mu_\plus u_\minus(j)]$.
Every $d\ge L_k/2$ has $d\ge n/2$, so $\nu_\plus(d)\le C_\nu r$. Every $d\in\mathcal V$
has $d\ge L_k/2$ and $|d-n|\le RB_N+1\le\delta n$, because $n\ge KB_N/(1-p)-1$, so
$|\nu_\plus(d)-\nu_\plus(n)|\le\varepsilon r$. Hence
\begin{align*}
 \Bigl|\Psi_N(k)-\nu_\plus(n)\bigl(1+\mu_\plus u_\minus(j)\bigr)\Bigr|
 &\le\sum_{d\notin\mathcal V}\nu_\plus(d)\Pi^\minus_{N-d}(j)
 +\sum_{d\in\mathcal V}\bigl|\nu_\plus(d)-\nu_\plus(n)\bigr|\Pi^\minus_{N-d}(j)\\
 &\quad+\nu_\plus(n)\,\varepsilon
 \ \le\ \bigl(C_\nu+2+\mu_\plus\bigr)\,\varepsilon\,r .
\end{align*}
Since $T_\plus(n)\le\widehat G_N(k)$, the decomposition gives the stated form, with the
double limit taken as before. At $k=N$ both $\Psi_N$ and the replacement vanish.
\end{proof}

\begin{proof}[Proof of Theorem~\ref{thm:preparation}(e)]
\emph{Reduction.} Fix $\varepsilon\in(0,\frac12)$. For large $N$ the set
$Z^\varepsilon_N$ lies in the zone of (c) for every $K$, because $\lambda_NB_N=o(N)$, and
it contains no $k=N$. So the decomposition of (c) gives
$\sup_{k\in Z^\varepsilon_N}|P_N(k)-G^\nu_N(k)-\Psi_N(k)|/\widehat G_N(k)\to0$, and
\eqref{eq:interior-target} holds if and only if $\Psi_N=o(\widehat G_N)$ uniformly on
$Z^\varepsilon_N$ for every $\varepsilon$. On $Z^\varepsilon_N$ we have
$T_\plus(N)\le\widehat G_N(k)\le\varepsilon^{-1}T_\plus(\lfloor\varepsilon N\rfloor)\le
C_\varepsilon T_\plus(N)$. So the condition reads
$\sup_{k\in Z^\varepsilon_N}\Psi_N(k)=o(T_\plus(N))$. Sufficiency of
$\bar\nu_\plus(n)=o(T_\plus(n))$ follows from~\eqref{eq:Psi-upper}.

For the necessity of the window condition, let $n$ be large, put $N=2n$ and choose
$k$ with $|L_k-n|\le1/(1-p)$, which is possible because $L_k$ moves in steps of
$1/(1-p)$. Then $k\in Z^{1/4}_N$. For $|d-n|\le B_n$ put $m=N-d$. Then
$m\in[N/4,3N/4]$, and $(m-j)-pm=(1-p)(L_k-d)$ has absolute value at most
$2(1-p)B_n\le c_{11}B_m$, with $c_{11}$ depending only on the laws, by
Lemma~\ref{lem:rv}(ii). Theorem~\ref{thm:bulk}(c) and Lemma~\ref{lem:profile}(d) give
$B_m\Pi^\minus_m(j)=B_mF^\minus_m(m-j)\ge c_{12}>0$ for large $n$. Also $d\ge L_k/2$ and
$B_m\le CB_n$. Hence $\Psi_N(k)\ge(c_{12}/C)W^{(N)}_n$ with $N=2n$. If~\eqref{eq:interior-target}
holds, then $\Psi_N=o(T_\plus(N))$ on $Z^{1/4}_N$ by the reduction, so
$W^{(2n)}_n=o(T_\plus(2n))=o(T_\plus(n))$.

\emph{Sufficiency of the window condition under (LR) for both signs.} Let
$k\in Z^\varepsilon_N$, with $j=N-k\ge(1-p)\varepsilon N$. A
nonzero term of~\eqref{eq:Psi} has $m=N-d\ge j$, so $m\in[(1-p)\varepsilon N,N]$ and
$d\in[\varepsilon N/2,N]$. Hence $B_m\ge c_\varepsilon B_N$, and
Theorem~\ref{thm:bulk}(b) gives $\Pi^\minus_m(j)\le C_\varepsilon/B_N$. Put
$\overline W_N=\sup_{\varepsilon N/4\le n\le N}W^{(N)}_n$, which is $o(T_\plus(N))$ by hypothesis,
and for a fixed law with $W_n=o(T_\plus(n))$ by (RV).
Split the $d$ axis into the blocks $\mathcal B_i=\{d:iB_N\le d-L_k<(i+1)B_N\}$,
$i\in\Z$. Each block meets $[\varepsilon N/2,N]$ in a set covered by at most
$C_\varepsilon$ windows $[n-B_n,n+B_n]$ with $\varepsilon N/4\le n\le N$, so
$\nu_\plus(\mathcal B_i\cap[\varepsilon N/2,N])\le C_\varepsilon B_N\overline W_N$. The
blocks with $|i|\le R$ contribute at most $(2R+1)C^2_\varepsilon\overline W_N$ to
$\Psi_N(k)$.

For the blocks with $|i|>R$ we do not use~\eqref{eq:unif-bigjump}, since its zone,
and its analogue at the minus end, carry logarithmic factors that are unbounded when
$\alpha_\plus\ne\alpha_\minus$. We use a bound that needs only an excess of order
$B_m$. Let $m$ be a large horizon and $k'>pm$ a point, and put $L=L^{(m)}_{k'}$ and
$g=\lceil L/32\rceil$. We apply~\eqref{eq:giant-split} at the horizon $m$ for a fresh
minus start and drop its last term. We bound its first term by
Lemma~\ref{lem:no-giant} with $\eta=1/32$ and $r=2$, and $\mathrm{Cen}_g$ and
$\mathrm{Int}_g$ by~\eqref{eq:giant-crude}, with $\kappa^{(\minus)}_\plus\le1$ and
$\gamma^{(\minus)}_\plus(m-k')\le(1+\mu_\plus)m$. By (LR) for the plus sign,
$f^*_\plus(g)\le C_{\mathrm{LR}}T_\plus(g)/g$, and $T_\plus(g)\le CT_\plus(L)$ by
Lemma~\ref{lem:rv}(ii). Since $L\le m$, this gives
\begin{equation}\label{eq:far-bound}
 F^\minus_m(k')\le C\,\frac{(mT_\plus(L))^2}{b_{\plus,m}}+\frac C{b_{\minus,m}}\,e^{-cL/b_{\minus,m}}
 +Ce^{-m^c}+C\,\frac mL\,T_\plus(L)\qquad(L\ge C_9B_m),
\end{equation}
with $C$ and $c$ depending only on the laws and $C_9\ge C_8$ the larger of the
thresholds of Lemma~\ref{lem:no-giant} at the two ends. The minus-end forms of
Lemmas~\ref{lem:giant} and~\ref{lem:no-giant}, with (LR) for the minus sign, give
the same bound for $k'<pm$ with the signs exchanged, with the same constants. There
$L$ is replaced by $L^\minus=(pm-k')/p$, $T_\plus$ by $T_\minus$, and $b_{\plus,m}$ and
$b_{\minus,m}$ by $b_{\minus,m}$ and $b_{\plus,m}$.

Now let $|i|>R$ and $d\in\mathcal B_i$, so that $m=N-d\in[(1-p)\varepsilon N,N]$ and
$\Pi^\minus_m(j)=F^\minus_m(k')$ with $k'=m-j$. Then $k'-pm=(1-p)(L_k-d)$. If $d<L_k$,
then $k'>pm$ and $L=L_k-d>(|i|-1)B_N$. If $d>L_k$, then $k'<pm$ and
$L^\minus=(1-p)(d-L_k)/p\ge(1-p)|i|B_N/p$. Put $c_*=\min(1,(1-p)/p)$, so the excess
is at least $y_i=c_*(|i|-1)B_N$ in both cases. Choose $R$ with $c_*R\ge\max(C_9,1)$.
Then $y_i\ge C_9B_N\ge C_9B_m$, so~\eqref{eq:far-bound} or its minus-end form
applies, and $y_i\ge B_N\ge b_{\sgn,N}$ for both signs. Each term
of~\eqref{eq:far-bound} decreases in $L$. For these $m$ we have $mT_\sgn\le NT_\sgn$
and $c_\varepsilon b_{\sgn,N}\le b_{\sgn,m}\le b_{\sgn,N}$. We treat the blocks with
$i<0$, and those with $i>0$ are the same with the signs exchanged. We multiply the
four terms of~\eqref{eq:far-bound} at $L=y_i$ by the block mass
$C_\varepsilon B_N\overline W_N$, sum over $i<-R$, and leave out the common factor
$C_\varepsilon\overline W_N$ below. Put $t=|i|-1\ge1$.

The first term. By~\eqref{eq:NT-bounds},
$(NT_\plus(y_i))^2\le16s^2(y_i/b_{\plus,N})^{-\alpha_\plus-1}$, so the term is at most
$C(B_N/b_{\plus,N})^{-\alpha_\plus}t^{-\alpha_\plus-1}\le Ct^{-\alpha_\plus-1}$. These sum
to a constant.

The second term. Put $x=B_N/b_{\minus,N}\ge1$. The term is at most $Cxe^{-cc_*tx}$,
and $\sum_{t\ge1}xe^{-cc_*tx}=x/(e^{cc_*x}-1)\le1/(cc_*)$.

The third term. At most $N/B_N+2$ blocks meet $[0,N]$, so these terms add up to at
most $C(N+2B_N)e^{-((1-p)\varepsilon N)^c}$, which tends to zero.

The fourth term. Fix $\delta'\in(0,\alpha_*)$. Since $c_*t\ge1$,
Lemma~\ref{lem:rv}(ii) gives $T_\plus(y_i)\le Ct^{-\alpha_\plus+\delta'}T_\plus(B_N)$. So
the term is at most $CNT_\plus(y_i)/t\le CNT_\plus(B_N)t^{-1-\alpha_\plus+\delta'}$, and
$NT_\plus(B_N)\le NT_\plus(b_{\plus,N})\to s$ by Lemma~\ref{lem:rv}(v). These sum to a
constant.

For $i>R$ the same four steps use~\eqref{eq:NT-bounds} for the minus sign,
$x=B_N/b_{\plus,N}$, and $NT_\minus(B_N)\le NT_\minus(b_{\minus,N})\to s$. So the blocks
with $|i|>R$ contribute at most $C_\varepsilon\overline W_N$ to $\Psi_N(k)$. Hence
$\sup_{k\in Z^\varepsilon_N}\Psi_N(k)\le C_{\varepsilon,R}\overline W_N=o(T_\plus(N))$.
\end{proof}

The gap between the two interior conditions is real. If a fixed law has $\nu_\plus(n_i)=T_\plus(n_i)$
on a sparse sequence and $\nu_\plus=0$ near it, then $W_n=O(T_\plus(n)/B_n)$, so under
(LR) for both signs the initial run is invisible on every interior zone, by the
window criterion (the sufficiency just proved). It is visible at the point $N-1$ for
$N=n_i+1$, where $P_N(N-1)\ge T_\plus(N-1)$, and $\nu_\plus(n)\ne o(T_\plus(n))$. So
the interior of the big-jump zone sees only window averages of the initial run,
while the edge of the zone sees single atoms.

\subsection{The stationary initial run}\label{sec:prep-stationary}

\begin{corollary}[stationary initial run]\label{cor:stationary}
Assume (RV) and (SPAN), and let the initial run be stationary,
$\nu_\sgn=T_\sgn/s$. Then the bulk criterion holds, and
\begin{enumerate}[(a)]
\item the atoms are $Q_\sgn(N)=s^{-1}\sum_{d\ge N}T_\sgn(d)\sim NT_\sgn(N)/(s(\alpha_\sgn-1))$.
\item if (LR) holds for the plus sign, then, with $\lambda_N$ as in~\eqref{eq:lambdaN},
\[
 \lim_{K\to\infty}\ \limsup_{N\to\infty}\ \sup_{pN+K\lambda_NB_N\le k<N}
 \Bigl|\frac{P_N(k)}{G^{\mathrm{st}}_N(k)}-1\Bigr|=0,
\]
where
$G^{\mathrm{st}}_N(k)=(N-L_k)f_\plus(\lfloor L_k\rfloor)/\mu_\minus+T_\plus(\lfloor L_k\rfloor)\,\kappa^{\mathrm{st}}_\plus(N-k)$,
and for fixed $v\in(p,1)$, with $\ell=(v-p)/(1-p)$,
$P_N(\lfloor vN\rfloor)/T_\plus(N)\to I_\plus(v)+\ell^{-\alpha_\plus}/\mu_\minus$. So in the
interior of the big-jump zone the kernel gains the term $T_\plus(\lfloor L_k\rfloor)/\mu_\minus$,
which is the complete initial plus run.
\item the boundary kernel at the plus end is
\[
 \frac{P_N(N-j)}{T_\plus(N)}\to\kappa^{\mathrm{st}}_\plus(j)=\frac{2}{s}\bigl(1+\mu_\plus u_\minus(j)\bigr)\qquad(j\ge1),
\]
while $P_N(N)/T_\plus(N)\to\infty$. Moreover $\kappa^{\mathrm{st}}_\plus(j)\to2/\mu_\minus$ as
$j\to\infty$, which is the value of $I_\plus(v)+\ell^{-\alpha_\plus}/\mu_\minus$ at $v=1$.
\end{enumerate}
\end{corollary}

\begin{proof}
The bulk criterion holds because $B_nT_\sgn(n)\to0$ by Lemma~\ref{lem:rv}(v). Part
(a) is Lemma~\ref{lem:rv}(iv).

We first compute $\kappa_\plus$. Here $\pi_\plus=\sum_dT_\plus(d)/s=\mu_\plus/s=p$, and
$\sum_{d=1}^jT_\minus(d)u_\minus(j-d)=1$ for $j\ge1$ by~\eqref{eq:cover}.
So~\eqref{eq:kappa-prep} gives $\kappa_\plus(j)=pu_\minus(j)+1/s$ for $j\ge1$.

(c) The function $\nu_\plus=T_\plus/s$ is regularly varying with index
$-\alpha_\plus<0$, so it satisfies (R$\nu$). Theorem~\ref{thm:preparation}(d) gives,
for $j\ge1$,
\[
 \frac{P_N(N-j)}{T_\plus(N)}\to pu_\minus(j)+\frac1s+\frac1s\bigl(1+\mu_\plus u_\minus(j)\bigr)
 =\frac2s+\frac{2\mu_\plus}{s}u_\minus(j),
\]
since $p=\mu_\plus/s$. Also $P_N(N)=Q_\plus(N)$, and $Q_\plus(N)/T_\plus(N)\to\infty$ by (a).
By~\eqref{eq:renewal-thm}, $\kappa^{\mathrm{st}}_\plus(j)\to(2/s)(1+\mu_\plus/\mu_\minus)=2/\mu_\minus$.

(b) By the regular case of Theorem~\ref{thm:preparation}(c), uniformly in the zone
with $k<N$ and in the double limit,
\[
 P_N(k)=G^\nu_N(k)+\frac{T_\plus(\lfloor L_k\rfloor)}{s}\bigl(1+\mu_\plus u_\minus(N-k)\bigr)
 +o\bigl(\widehat G_N(k)\bigr),
\]
where we used $T_\plus(\lfloor L_k\rfloor)\le\widehat G_N(k)$. By the computation of
$\kappa_\plus$ the right side is $G^{\mathrm{st}}_N(k)+o(\widehat G_N(k))$. For the ratio
form we need $G^{\mathrm{st}}_N(k)\ge c\widehat G_N(k)$. Since $\kappa^{\mathrm{st}}_\plus\ge2/s$,
this holds when $L_k>N/2$. When $L_k\le N/2$, (LR) gives
$(N-L_k)f_\plus(\lfloor L_k\rfloor)/\mu_\minus\ge c(N/L_k)T_\plus(\lfloor L_k\rfloor)$ for large
$N$. Now fix $v\in(p,1)$ and let $k=\lfloor vN\rfloor$. Then $k$ lies in the zone for
every $K$ when $N$ is large, because $\lambda_NB_N=o(N)$. Also $L_k=\ell N+O(1)$
and $N-k\to\infty$. By (LR) and (RV),
$(N-L_k)f_\plus(\lfloor L_k\rfloor)/T_\plus(N)\to(1-\ell)\alpha_\plus\ell^{-\alpha_\plus-1}$ and
$T_\plus(\lfloor L_k\rfloor)/T_\plus(N)\to\ell^{-\alpha_\plus}$, and
$\kappa^{\mathrm{st}}_\plus(N-k)\to2/\mu_\minus$. This gives the limit
$I_\plus(v)+\ell^{-\alpha_\plus}/\mu_\minus$, with $I_\plus$ as in
Corollary~\ref{cor:infinite-density}. At $v=1$ it equals $2/\mu_\minus$.
\end{proof}

In $\kappa^{\mathrm{st}}_\plus(j)$ the term $pu_\minus(j)$ comes from a short first plus
run, the term $1/s$ from a short first minus run, and the term $(1+\mu_\plus u_\minus(j))/s$
from a first plus run that ends within $O(1)$ steps of the horizon.
Section~\ref{sec:verify} reports a check of part (c) against the exact law.
Wang, Vezzani, Burioni and Barkai~\cite[Eqs.~(28), (46) and~(47)]{wvbb2019}
compared the equilibrated and the ordinary start formally, for a biased
continuous-time random walk. The equilibrated start adds the term $\xi^{-\alpha}$
to their infinite density, where $\xi$ is the fraction of the time taken by the
longest trap, and an atom at the origin of order $t^{1-\alpha}$. They correspond
to the term $\ell^{-\alpha_\plus}/\mu_\minus$ of (b) and to the atom of (a).

The same index comparisons place other initial runs. If
$Q_\sgn\in\RV_{-\rho}$ with $\rho>0$ and $\nu_\sgn(n)\sim\rho Q_\sgn(n)/n$, the bulk
criterion holds, the criterion for the whole range holds for $\rho>1/\alpha_*$ and
fails for $\rho<1/\alpha_*$, and since $\nu_\plus/T_\plus\in\RV_{\alpha_\plus-1-\rho}$
the initial run is invisible in the boundary layer and, under (LR) for the plus
sign, in the big-jump zone for $\rho>\alpha_\plus-1$ and visible for
$\rho<\alpha_\plus-1$. The stationary initial run sits on the line
$\rho=\alpha_\sgn-1$ with $\nu_\sgn/T_\sgn=1/s$, so it is visible there.

\subsection{No universal trim\ifarxivbuild{} and identification\fi}\label{sec:prep-trim}

Part (b) of Theorem~\ref{thm:preparation} holds on every macroscopic interior with
no condition. No sublinear trim of the range does the same.

\begin{corollary}[no universal trim]\label{cor:no-trim}
Assume (RV) and (SPAN). For every sequence of nonnegative integers $r_N=o(N)$
there is an initial run with a fixed deterministic sign for which the bulk
approximation fails on $r_N\le k\le N-r_N$ along a subsequence.
\end{corollary}

\begin{proof}
Put $R(n)=\max_{1\le i\le2n}r_i$ and
\[
 m(n)=\Bigl\lceil\max\Bigl\{\sqrt n,\ \frac{4(R(n)+1)}{1-p}\Bigr\}\Bigr\rceil .
\]
Since $r_i=o(i)$, for every $\eta>0$ we have $r_i\le\eta i$ for $i\ge i_\eta$, so
$R(n)\le\max_{i<i_\eta}r_i+2\eta n$ and $R(n)=o(n)$. Hence $m(n)\to\infty$ and
$m(n)=o(n)$. Since $B$ is regularly varying with the positive index $1/\alpha_*$,
Lemma~\ref{lem:rv}(ii) gives $B_{m(n)}/B_n\to0$. Choose $n_1<n_2<\cdots$ so that
$n_{i+1}>2n_i$, $m(n_i)\le n_i$ and $\sum_iB_{m(n_i)}/B_{n_i}\le\frac12$. Put
$m_i=m(n_i)$, $N_i=n_i+m_i\le2n_i$ and $v_i=B_{m_i}/B_{n_i}$. The initial run has
$\Sigma=\plus$ almost surely, $\nu_\plus(n_i)=v_i$ for every $i$, and the remaining mass
at $d=1$. Then $D$ is finite almost surely.

Fix $i$ and let $k'_i=\lfloor pm_i\rfloor$. By Theorem~\ref{thm:bulk}(c) and
Lemma~\ref{lem:profile}(d) there is $c>0$ with $B_{m_i}F^\minus_{m_i}(k'_i)\ge c$ for
large $i$. Let $k_i=n_i+k'_i$. Proposition~\ref{prop:first-run}, keeping only the term
$d=n_i$, gives
\[
 B_{N_i}P_{N_i}(k_i)\ge B_{N_i}\,v_i\,F^\minus_{m_i}(k'_i)\ge c\,\frac{B_{N_i}}{B_{n_i}}\ge c,
\]
because $B$ is nondecreasing. On the other hand
$k_i-pN_i=(1-p)n_i+(k'_i-pm_i)\ge(1-p)n_i-1$, so
$x_{N_i,k_i}\ge((1-p)n_i-1)/B_{2n_i}\to\infty$, and
$h_{N_i}(x_{N_i,k_i})\to0$ by Lemma~\ref{lem:profile}(b). So
$|B_{N_i}P_{N_i}(k_i)-h_{N_i}(x_{N_i,k_i})|\ge c/2$ for large $i$.

It remains to check that $k_i$ is retained. Since $N_i\le2n_i$, we have
$r_{N_i}\le R(n_i)=o(n_i)$, so $k_i\ge n_i\ge r_{N_i}$ for large $i$. Also
$N_i-k_i=m_i-k'_i\ge(1-p)m_i\ge4(R(n_i)+1)>r_{N_i}$. Hence
$r_{N_i}\le k_i\le N_i-r_{N_i}$, and the approximation fails along $N_i$.
\end{proof}

The initial run violates the bulk criterion, since
$B_{n_i}\nu_\plus(n_i)=B_{m_i}\to\infty$.
\ifarxivbuild
Finally, the boundary layer is a renewal mass
function, which determines its interarrival law, so the ends identify the run laws.

\begin{remark}[identification from the ends]\label{thm:inverse}\label{rem:thinning}
Assume (RV) for both signs.
(a) For a fresh plus start $P_N(N)=T_\plus(N)$ by
Proposition~\ref{prop:fresh-identity}, so the atoms determine $f_\plus$. By
Theorem~\ref{thm:bigjump}(a) and~\eqref{eq:kappa}, $P_N(N-j)/P_N(N)\to u_\minus(j)$ for
every fixed $j\ge1$. Splitting the sum that defines $u_\minus(j)$ according to the
length of the first run gives the renewal equation
\begin{equation}\label{eq:renewal-inverse}
 f_\minus(j)=u_\minus(j)-\sum_{i=1}^{j-1}f_\minus(i)\,u_\minus(j-i),\qquad j\ge1,
\end{equation}
so $u_\minus$ determines $f_\minus$ by induction on $j$, and the family $(P_N)_{N\ge1}$
determines both run laws.
(b) Let the initial run satisfy $Q_\plus(N)/T_\plus(N)\to q\in(0,\infty)$. Then
$\pi_\plus=\Prob(\Sigma=\plus)>0$ and $\nu_\plus(n)=Q_\plus(n)-Q_\plus(n+1)=o(T_\plus(n))$,
since $T_\plus(n+1)/T_\plus(n)\to1$. Theorem~\ref{thm:preparation}(d) and
$P_N(N)=Q_\plus(N)$ give $P_N(N-j)/P_N(N)\to\kappa_\plus(j)/q$. If the law of the
initial run and $q$ are known, then
$\pi_\plus u_\minus(j)=\kappa_\plus(j)-\sum_{d=1}^{j}\nu_\minus(d)u_\minus(j-d)$ determines
$u_\minus$ by induction, and~\eqref{eq:renewal-inverse} gives $f_\minus$. The stationary
initial run fails the hypothesis, since its atom is of order $NT_\plus(N)$ by
Corollary~\ref{cor:stationary}(a).
(c) Ratios at the minus end cannot replace the atoms. For a fresh plus start
Theorem~\ref{thm:bigjump}(a) with the signs exchanged gives
$P_N(j)/T_\minus(N)\to u_\plus(j)$, so the ratios $P_N(j)/P_N(i)$ determine $u_\plus$
only up to a positive factor. For $0<c\le1$ let $X^c_\plus$ be the sum of a geometric
number $G$ of independent copies of $X_\plus$, with $\Prob(G=n)=c(1-c)^{n-1}$. With
generating functions $F$ of $f_\plus$ and $U=1/(1-F)$, the law of $X^c_\plus$ has
$F_c=cF/(1-(1-c)F)$ and $1/(1-F_c)=1+c(U-1)$, so its renewal mass function is
$cu_\plus$ on $j\ge1$. Its support contains that of $X_\plus$, which gives (SPAN). Its
tail is $\Prob(X^c_\plus\ge n)\sim T_\plus(n)/c$, by $\Prob(S^\plus_k\ge n)\sim kT_\plus(n)$
for each $k$~\cite[Theorem~2.1]{berger2019ptrf}, Kesten's
bound~\cite[Lemma~1.3.5]{ekm1997} and dominated convergence, which gives (RV).
Replacing $X_\plus$ by $X^c_\plus$ changes no ratio at either end. The atoms separate
the laws $f^c_\plus$, and so does the bulk, since the mean $\mu_\plus/c$ changes $p$.
\end{remark}
\else
Finally, the boundary layer is a renewal mass function, which determines its interarrival
law, so the ends identify both run laws. The arXiv version of this paper makes this precise
in a remark at the end of this section, for fresh and prepared starts, and shows by a
geometric thinning of the runs that the ratios at the ends alone do not.
\fi
